% Part: normal-modal-logic
% Chapter: axioms-systems
% Section: introduction

\documentclass[../../../include/open-logic-section]{subfiles}

\begin{document}

\olfileid[hi]{nml}{axs}{int}

\olsection{परिचय}

हमारे पास मोडल प्रतिरूपों के आधार पर मूल मोडल भाषा का अर्थविज्ञान और किसी
!!{formula} के मान्य होने की धारणा है---सभी प्रतिरूपों के सभी जगतों पर सत्य---
या प्रतिरूपों अथवा फ़्रेमों के किसी वर्ग के सापेक्ष मान्य होने की धारणा---
वर्ग के सभी प्रतिरूपों, या फ़्रेम पर आधारित सभी प्रतिरूपों, के सभी जगतों पर
सत्य। तर्कशास्त्र सामान्यतः मान्यता के ऐसे अर्थगत अभिलक्षणों को
!!{derivability} की प्रमाण-सैद्धांतिक धारणा से जोड़ता है। लक्ष्य किसी तंत्र
में !!{derivability} की ऐसी धारणा परिभाषित करना है कि कोई !!{formula} तभी
!!{derivable} हो जब वह मान्य हो।

सबसे सरल और ऐतिहासिक रूप से सबसे पुराने !!{derivation}-तंत्र तथाकथित
हिल्बर्ट-प्रकार या अभिगृहीतीय !!{derivation}-तंत्र हैं। अनेक मोडल
तर्कशास्त्रों के लिए हिल्बर्ट-प्रकार !!{derivation}-तंत्र बनाना अपेक्षाकृत
आसान है: अधिसैद्धांतिक अध्ययन की वस्तुओं के रूप में वे सरल हैं (जैसे
यथार्थता और पूर्णता सिद्ध करने के लिए)। किंतु उनमें !!{formula} सिद्ध करना,
उदाहरणार्थ प्राकृतिक निगमन तंत्रों की तुलना में, कहीं कठिन है।

हिल्बर्ट-प्रकार !!{derivation}-तंत्र में किसी !!{formula} की !!{derivation},
कुछ अभिगृहीतों से कुछ अनुमान-नियमों के माध्यम से उस !!{formula} तक ले जाने
वाला सूत्रों का अनुक्रम है। चूँकि हम चाहते हैं कि !!{derivation}-तंत्र
अर्थविज्ञान से मेल खाए, हमें सुनिश्चित करना है कि !!{derivable} सूत्रों का
समुच्चय सभी प्रतिरूपों में सत्य हो (या उन सभी प्रतिरूपों में सत्य हो जिनमें
सभी अभिगृहीत सत्य हैं)। पहले हम मोडल तर्कशास्त्रों के वे गुण अलग करेंगे जो
इसके लिए आवश्यक हैं: ``सामान्य'' मोडल तर्कशास्त्र। सामान्य मोडल तर्कशास्त्र
के लिए केवल दो अनुमान-नियम मानने होते हैं: मोडस पोनेन्स और आवश्यकीकरण।
अभिगृहीतों के रूप में हम सर्वसत्यताओं के सभी प्रतिस्थापन-दृष्टांत और संबंधित
मोडल तर्कशास्त्र के अनुसार कुछ मोडल अभिगृहीत लेते हैं। केवल सभी प्रतिरूपों
के वर्ग में रुचि होने पर भी हमें $\Ax{K}$\iftag{notprvDiamond}{}{ और~$\Ax{Dual}$}
के सभी प्रतिस्थापन-दृष्टांत अभिगृहीत मानने होते हैं। इतने से ही न्यूनतम
सामान्य मोडल तर्कशास्त्र~$\Log K$ उत्पन्न होता है।

\begin{defn}
\emph{मोडस पोनेन्स} का नियम यह अनुमान-स्कीमा है:
\begin{prooftree}
\AxiomC{$!A$}
\AxiomC{$!A \lif !B$}
\RightLabel{\MP}
\BinaryInfC{$!B$}
\end{prooftree}
हम कहते हैं कि !!{formula}~$!B$, मोडस पोनेन्स से !!{formula}ों $!A$, $!C$ से
\emph{निकलता है}, तभी जब $!C \ident !A \lif !B$।
\end{defn}

\begin{defn}
\emph{आवश्यकीकरण} का नियम यह अनुमान-स्कीमा है:
\begin{prooftree}
\AxiomC{$!A$}
\RightLabel{\Nec}
\UnaryInfC{$\Box !A$}
\end{prooftree}
हम कहते हैं कि !!{formula}~$!B$, आवश्यकीकरण से !!{formula}~$!A$ से निकलता है,
तभी जब $!B \ident \Box !A$।
\end{defn}

\begin{defn}
अभिगृहीतों के समुच्चय~$\Sigma$ से \emph{!!{derivation}}, !!{formula}ों
$!B_1$, $!B_2$, \dots, $!B_n$ का ऐसा अनुक्रम है जिसमें प्रत्येक $!B_i$
निम्न में से कोई है:
\begin{enumerate}
\item किसी सर्वसत्यता का प्रतिस्थापन-दृष्टांत; या
\item $\Sigma$ के किसी !!{formula} का प्रतिस्थापन-दृष्टांत; या
\item $j,k<i$ वाले दो !!{formula}ों $!B_j$, $!B_k$ से मोडस पोनेन्स द्वारा
  निकलता है; या
\item $j<i$ वाले किसी !!{formula}~$!B_j$ से आवश्यकीकरण द्वारा निकलता है।
\end{enumerate}
यदि $!B_n \ident !A$ वाली ऐसी कोई !!{derivation} हो, तो हम कहते हैं कि
$!A$, $\Sigma$ से \emph{!!{derivable}} है; संकेत में
$\Sigma \Proves !A$।
\end{defn}

इस परिभाषा से यह निकलेगा कि !!{derivable} सूत्रों का समुच्चय सामान्य मोडल
तर्कशास्त्र बनाता है और प्रत्येक !!{derivable} !!{formula} उस प्रत्येक
प्रतिरूप में सत्य है जिसमें प्रत्येक अभिगृहीत सत्य हो। !!{derivation}यों का
यह गुण \emph{यथार्थता} कहलाता है। इसका विलोम, \emph{पूर्णता}, सिद्ध करना
अधिक कठिन है।

\end{document}
